\section{Limitations}
\label{sec:limitations}
\label{page:limitations_start}

\subsection{Scale, stated as a falsifiable prediction}

The study operates at 5.6M parameters, $d_{\text{model}} = 256$, one block, a
vocabulary of 8{,}192, and three epochs of WikiText-103. Sparse-MoE benefits are
known to emerge with scale, so a reviewer can reasonably say the negative result
may be an artefact of operating below the regime where sparsity pays. \textbf{That
objection is correct and cannot be answered with the runs that exist.}

Rather than hedge it, we convert it into the paper's forward claim. The account in
\S\ref{sec:dilution} is that routing divides the \emph{diversity} of data each
parameter block sees, not its update budget---which the arithmetic there shows is
matched exactly. That account predicts a specific crossover: as the corpus grows at
fixed parameters, each expert's router-selected slice comes to contain enough of
the distribution's shared structure that the re-learning penalty amortises, so
\textbf{the MoR$-$MoRE gap should narrow monotonically with training data and the
ordering should eventually invert in favour of the composition}. The out-of-domain
result (\S\ref{sec:ood}) is weak early evidence in this direction, since MoRE is
already the better arm on data it was not trained on.

The measurement that would falsify it: hold parameters, depth budget and optimiser
fixed; sweep training tokens over at least a decade; and report the MoR$-$MoRE gap
as a function of tokens per expert. If the gap is flat or widening, the diversity
account is wrong and the composition's deficit is structural. We did not run that
sweep, and a second model size---the single most valuable follow-up---is likewise
unrun.

% >>>ACLMOVE:budget|\subsection{Budget, not convergence}\label{sec:budget} The 3-epoch budget leaves every arm short of convergence (\S\ref{sec:protocol}), and the arms are furthest apart on that axis at the cutoff: the final-epoch drop is $0.091$ nats for MoE, $0.132$ for MoR and $0.176$ for MoRE, so \textbf{the arm that loses in domain is the one still improving fastest}. The ordering is budget-dependent in the direction this paper's mechanism predicts---a routing arm that re-learns shared structure more slowly (\S\ref{sec:dilution}) should be behind at a fixed budget and closing---and we do not extrapolate a crossing epoch, because the curves decelerate. Appendix~\ref{app:budget} gives the per-epoch curves and the rest of this section.
\subsection{Budget, not convergence}
\label{sec:budget}

Three epochs is a compute budget, and no arm is trained to convergence
(\S\ref{sec:protocol}). This interacts badly with exactly the mechanism the paper
studies: \S\ref{sec:saturation} shows depth allocation is demonstrably
budget-dependent, so the behaviour at 10 epochs is genuinely unknown and might
differ in the direction that matters. Every arm receives the identical budget, so
the comparison is internally valid; the generalisation to trained-to-convergence
models is not established.

\textbf{The arms are not equally far from convergence, and the ordering tracks
that gap.} The admitted runs log validation loss once per epoch, so the curve over
the three epochs is available for all fifteen cells. The final-epoch drop is
$0.091$ nats for MoE, $0.132$ for MoR and $0.176$ for MoRE: \textbf{the arm that
loses in domain is the one still improving fastest when the budget ends.} We state
the differential and not a crossing point---the curves decelerate, so extrapolating
an epoch at which MoRE overtakes MoR would be inventing a number---but it makes the
in-domain ordering explicitly budget-dependent, in the direction this paper's
mechanism predicts: \S\ref{sec:dilution} argues a routing arm re-learns shared
structure more slowly, which implies it should be behind at a fixed budget and
closing.

\begin{center}
{\dispsize
\setlength{\tabcolsep}{9pt}\renewcommand{\arraystretch}{1.25}
\small
\begin{tabular}{lccc}
\toprule
validation task loss & MoE & MoR & MoRE \\
\midrule
epoch 1 & 4.1453 & 3.8202 & 3.9416 \\
epoch 2 & 4.0392 & 3.6472 & 3.8131 \\
epoch 3 (reported) & 3.9483 & 3.5154 & 3.6370 \\
\midrule
final-epoch drop & $+0.0909$ & $+0.1318$ & $\mathbf{+0.1761}$ \\
\bottomrule
\end{tabular}
}
\end{center}

The curves also show a second quantity moving in the same direction. MoRE's
between-seed standard deviation tightens across the three epochs
($0.0338 \rightarrow 0.0211 \rightarrow 0.0134$) while MoR's stays between
$0.033$ and $0.039$. That is direct support for the parsimonious reading of the
variance-ratio result in \S\ref{sec:quality}---a policy collapsing to a
near-constant has less left to vary---which we otherwise report as open.

Relatedly, one hyperparameter configuration per arm was frozen from a grid run on a
corpus roughly 60$\times$ smaller, and \S\ref{sec:saturation} demonstrates that at
least one weight---the ponder penalty---did not transfer. We did not re-tune per
arm, so the headline ordering could in principle be re-ordered by a per-architecture
search. Given the gaps involved (0.12 nats between the recursive arms), we regard
that as a real risk for the MoR--MoRE comparison and a small one for the comparisons
against MoE.
% <<<ACLMOVE:budget

\subsection{Statistical resolution}

$n = 5$ seeds puts every headline comparison exactly on the design's resolution
floor of $p = 0.00794$ (\S\ref{sec:stats}). ``As extreme as a five-seed design
admits'' is the strongest statement available, and it is weaker than significance.
No larger effect could produce a smaller $p$ here; only more seeds could. Cohen's
$d$ values are reported alongside every gap precisely because a large $d$ on a
small gap is the misreading to avoid---across our own re-exports, $d$ has moved by
a factor of 1.5 while the underlying gap moved by less than 0.005 nats.

\subsection{What was declared and not run}
\label{sec:declared}

One comparison in the pre-registered Family A, a depth-matched fixed-depth control,
was declared before results existed and \emph{still consumes its Holm slot}: we do
not delete it to loosen the thresholds on the survivors. The forced-depth sweep in
\S\ref{sec:halting} answers the same question using inference passes over the
canonical checkpoints rather than a trained control arm, so we report it as an
exploratory mechanism result and leave the declared slot consumed and unfilled. A
fixed-depth sweep trained from scratch, a shuffled-depth control, and a
depth--input mutual-information analysis were all declared exploratory and are
unrun.

\subsection{Evaluation caveats}

BLiMP and EWoK are scored with the tokenizer fitted to our training corpus
(8{,}192 BPE merges), so \textbf{absolute accuracies are not comparable to
published models} using 32k--50k vocabularies; only the cross-arm comparison is
valid. EWoK sits within noise of its 50\% chance floor for every arm, which we
report as a scale limitation rather than a result---at 5.6M parameters a world
knowledge probe is below its useful range. All evaluations outside the two declared
families are uncorrected and carry no significance verdict in the abstract or
conclusions.

The perplexities in Table~\ref{tab:main} are intra-study quantities for the same
reason and should not be compared to published WikiText-103 numbers.

\subsection{Disclosures}
\label{sec:honesty}

Several facts about this study would count against it if a reader discovered them
independently, so we state them.

\textbf{An earlier version of the composed arm was not the architecture we define,
and fixing it made our headline worse.} It re-routed at every recursion step,
letting a token traverse up to seven \emph{different} expert FFNs---a heterogeneous
seven-layer network, not one block iterated, and strictly more expressive than the
specification allows. That arm scored 3.5031 and was statistically indistinguishable
from MoR. The architecturally-correct route-once arm reported here scores 3.6370,
\textbf{0.134 nats worse}. The earlier arm's apparent parity with MoR was an
artefact of a specification violation. We retain it unmodified as a labelled
ablation.

\textbf{The pre-registration is partial and the registration file says so itself.}
MoE and MoR were already trained when the test families were declared, so only the
comparisons involving the composed arm are genuinely pre-result. The MoR-versus-MoE
comparison was specified after its data existed, and we do not describe the
registration as complete.

\textbf{The matrix spans two GPUs and four code revisions.} Two arms were
retrained on one card after the original checkpoints were lost; the composed arm
remains on the original hardware, and comparability rests on a frozen
specification and a seed-blind configuration-equality check rather than on a
single revision. We have direct evidence the hardware split is benign: both
retrained arms reproduced their published means to $-0.117\sigma$ and
$-0.119\sigma$---agreement to the third decimal across arms sharing no
parameters, a \emph{common-mode} offset rather than sampling noise, and one that
cancels in the between-arm gaps carrying every verdict (the MoE$-$MoR gap moved
by 0.000185 nats). Per-seed values do \emph{not} reproduce and we do not claim
they do; the reproducible quantity is the arm mean.

\textbf{We report no per-token depth-allocation error.} English supplies no
ground-truth depth per token, so that metric is \texttt{N/A}---not zero. A zero
would read as perfect allocation against a curriculum that does not exist.

\section{Conclusion}
\label{sec:conclusion}

At 5.6M parameters on WikiText-103, composing sparse expert routing with
adaptive recursion buys efficiency and robustness rather than peak quality. All
three arms lie on the quality/compute frontier; MoRE buys the most quality per
unit of compute ($72\%$ of recursion's gain for $24\%$ of its compute) and out of
domain \emph{dominates} recursion alone, reaching lower loss at a quarter of the
compute---the one respect in which composing the mechanisms does something
neither does alone. What it does not buy is in-domain quality: recursion alone
leads by 0.1216 nats, significant at our resolution floor. The mechanism is not
the obvious one---routing does not dilute the update budget, and subspace
trapping fails three ways---but a competition for one per-token budget, in which
the arm preserving specialisation saturates its depth. We release the engine, the
evaluation battery, and the forced-depth control, and it is the last we would
most like to see reused: it costs one inference pass, and it is what separates
two arms that are indistinguishable by mean-and-histogram.
